\section{Historical Cohorts and Era Boundaries}
\label{app:history}
Git history contains completed experiments from earlier implementations, as well as the current deployment records. We bind each recovered cohort to its recorded model, source revision, and active mechanisms. A successful earlier test can describe that implementation without qualifying its successor. In particular, the July FP8 tree, the September Cat10 replay study, and the deployed NVFP4 Hydra27 route are distinct systems. The Cat10 study is superseded even though it ran after the August deployment.

\paragraph{Superseded FP8 application cohort}
The July 22 archive contains one attempt per arm on the same sixteen SWE-bench Verified Astropy tasks: a branching head with a six-token suffix tail, native MTP-5, and native MTP-11. All three boot logs identify the local model alias \texttt{qwen3.6-27b-fp8}, FP8 quantization, and a four-sequence server. The archive records source revisions \texttt{224c5db}, \texttt{103cedc}, and \texttt{2fc0faf}, respectively. Prefix caching is enabled; the tree receipt additionally records recurrent replay and the KV/slot remapping flags. These model aliases and source records do not authenticate an immutable checkpoint checksum or loaded-binary identity. They do establish that this is an older FP8 cohort, rather than a native control for the NVFP4 deployments in Section~\ref{sec:agent-workload}~\cite{lumo2026archive}.

\begin{table}[t]
\caption{Superseded July 22 FP8 cohort: one attempt per arm on the same sixteen Astropy tasks. Counts come from all 48 saved evaluator reports. Summed agent time is accumulated task time under a four-sequence server, not elapsed campaign time or isolated decoding time.}
\label{tab:historical-fp8}
\centering\small
\begin{tabular}{@{}lrrr@{}}
\toprule
Historical arm & Attempts & Resolved & Agent min \\
\midrule
Tree tail6 & 16 & 10 & 611.20 \\
Native MTP-5 & 16 & 10 & 373.73 \\
Native MTP-11 & 16 & 8 & 557.12 \\
\bottomrule
\end{tabular}
\end{table}


Recovery from Git commit \texttt{51b7dafdf} reconciles all 48 evaluator reports with their corresponding runner metadata and all three campaign summaries. No adverse attempt is removed. The tree has five failed-test patches and one empty-patch/application failure; MTP-5 has six failed-test patches, and MTP-11 has eight. None of these runner records reports a timeout. Equal aggregate resolve counts do not imply equal behavior: tree and MTP-5 resolve nine tasks in common, each resolves one task the other does not, and both fail the remaining five. Tree and MTP-11 share eight resolved tasks; tree alone resolves two additional tasks. With one stochastic attempt per arm and task, these are recorded outcomes, not evidence of quality equivalence or a benchmark-wide success rate.

The original closeout's decode proxies, including 32.85, 42.74, and 39.95 tokens/s, remain excluded. Its global accepted-token counters and retained pure-decode wall intervals cover different event populations; the subtraction residual is not measured host cost. Task-level agent durations survive that measurement correction, but they include tool work and differing trajectories. Table~\ref{tab:historical-fp8} therefore reports their accumulated time without interpreting it as a current decoder speedup.

\paragraph{Other recovered comparator coverage}
Four archived SGLang calibration records contain 56 completed synthetic requests in total: two EAGLE controls at maximum concurrency one, and DSpark at maximum concurrency one and eight. Their input and output lengths are 1,024 tokens. The safe projections retain original Git blob identities, model paths, algorithm settings, seeds, token totals, and benchmark durations; recomputing output tokens divided by duration reproduces each recorded aggregate output rate. These calibration measurements include first-output and concurrency effects and use different random seeds. They are neither agent-task rates under Eq.~\ref{eq:pooled-decode} nor a matched comparison with the current tree verifier~\cite{lumo2026archive}.

\paragraph{Qualification evidence does not transfer by name}
Two July 31 records provide narrower primary evidence. A shadow qrow16 attention check at one layer and sequence length 22,901 records zero mismatches over 393,216 output bytes and 3,072 log-sum-exp bytes. Its reference-served run finalized 810 decode events with no pending work. A separate FP8-era B4 batched-GDN check records zero byte mismatches in 432 candidate-surface and 288 graph-baseline comparisons across 48 layers. Reference state was restored and served, and the campaign later failed its snapshot/finalization boundary. That component result does not turn the rejected campaign into a valid lifecycle or performance result.

Earlier replay and fixed32 credentials refer to their original kernels and protocols. The older shared-step replay implementation differs from Hydra27's two-level scan and native per-layer committer. An August 15 widened gate explicitly invalidated earlier source-v7 fixed32 credentials. Such records remain historical, including the original failures and subsequent corrections; they do not override the current native-reference result in Section~\ref{sec:current-validation}.

The current NVFP4 Cqc10 receipt closes 49,333 recorded work-census events, with no pending verification, drafting, or commitment intervals. This is positive terminal bookkeeping evidence for that deployment, not a recurrent-state or next-forward-logit oracle. The source-bound recovery ledger distinguishes primary records reproduced here, previously integrated current-route evidence, superseded credentials, and reported results whose raw payload was not recovered. No new inference was needed to recover these observations.
